%==============================================================================
%  MSDS 498 Capstone --- Test Plan
%  AI-Driven High School Data Intelligence Platform
%  Typeset to match the report style of MSDS_498_version-wk3.pdf
%  Compile with: pdflatex MSDS_498_Test_Plan.tex  (run twice for the TOC)
%==============================================================================
\documentclass[12pt,letterpaper]{article}

\usepackage[T1]{fontenc}
\usepackage[utf8]{inputenc}
\usepackage[margin=1in]{geometry}
\usepackage[expansion=false]{microtype}
\usepackage{setspace}
\usepackage{booktabs}
\usepackage{longtable}
\usepackage{array}
\usepackage{ragged2e}
\usepackage{tabularx}
\usepackage{listings}
\usepackage{xcolor}
\usepackage{caption}
\usepackage{textcomp}
\usepackage[hidelinks]{hyperref}
\usepackage{underscore} % lets snake_case identifiers be typed literally

%------------------------------------------------------------------ spacing ---
\setstretch{1.15}
\setlength{\parindent}{1.5em}
\captionsetup{font=small,labelfont=bf,skip=6pt}
\captionsetup[longtable]{skip=4pt}

%--------------------------------------------------------------- convenience ---
\newcommand{\ra}{\textrightarrow{}}

% \code{...}: typewriter text that also permits line breaks after "/" and "."
% (the underscore package already makes "_" a breakpoint). Re-scans its argument
% with those two characters active, so long paths/filenames wrap inside table
% cells instead of overflowing the margin. Escaped %, {, }, [, ], and math all
% survive the re-scan unchanged.
\begingroup
\catcode`\/=\active \catcode`\.=\active
\gdef/{\char`\/\allowbreak}
\gdef.{\char`\.\allowbreak}
\endgroup
\newcommand{\code}[1]{{\ttfamily\endlinechar=-1
  \catcode`\/=\active \catcode`\.=\active \scantokens{#1}}}

% Gentle line-breaking tolerance so justified prose doesn't poke into the margin.
\tolerance=1200
\emergencystretch=2em

%----------------------------------------------------------------- listings ----
\definecolor{codebg}{gray}{0.965}
\definecolor{codeframe}{gray}{0.80}
\definecolor{kw}{rgb}{0.10,0.10,0.55}
\definecolor{cmt}{rgb}{0.30,0.45,0.30}
\definecolor{str}{rgb}{0.55,0.12,0.12}

\lstdefinestyle{shell}{
  basicstyle=\ttfamily\footnotesize,
  backgroundcolor=\color{codebg},
  frame=single, framesep=6pt, rulecolor=\color{codeframe},
  xleftmargin=8pt, xrightmargin=8pt,
  breaklines=true, breakatwhitespace=false,
  columns=fullflexible, keepspaces=true,
  showstringspaces=false, aboveskip=1.1em, belowskip=1.1em,
}
\lstdefinestyle{sql}{
  style=shell,
  language=SQL,
  keywordstyle=\color{kw}\bfseries,
  commentstyle=\color{cmt}\itshape,
  stringstyle=\color{str},
  morekeywords={ROUND,AVG,COUNT},
}

% Uniform header cells for the tables
\newcommand{\hd}[1]{\textbf{#1}}

\hypersetup{
  pdftitle={MSDS 498 Capstone Test Plan},
  pdfauthor={Max Chalekson, Sheng Hu, Qifan Yang},
}

%==============================================================================
\begin{document}

%---------------------------------------------------------------- title page ---
\begin{titlepage}
\centering
\thispagestyle{empty}
\vspace*{1.6in}

{\LARGE\bfseries Test Plan\par}
\vspace{0.9em}
{\large Unit, Integration, System, and User Acceptance Testing\\[0.3em]
for the AI-Driven High School Data Intelligence Platform\par}

\vspace{1.4in}

{\large Max Chalekson\\ Sheng Hu\\ Qifan Yang\par}

\vspace{0.8in}

{\normalsize Northwestern University\\ Master of Science in Data Science\par}

\vspace{0.7in}

{\normalsize Prepared for Bob Henkins\\ Northwestern Undergraduate Admissions and Enrollment\par}

\vfill
{\normalsize Summer 2026\par}
\end{titlepage}

%------------------------------------------------------------- table of cont ---
\clearpage
\pagenumbering{roman}
\setcounter{tocdepth}{1}
\tableofcontents
\clearpage
\pagenumbering{arabic}

%------------------------------------------------------------------- overview ---
\section*{Overview}
\addcontentsline{toc}{section}{Overview}

This plan covers unit, integration, system, and user acceptance testing for the
Northwestern high-school data platform (Sections 3--4 of
\code{written-report-iterations/MSDS_498_version-wk3.pdf}). It is not a generic
template --- every scenario below is grounded in a specific, real risk already
found in this pipeline, several of them actual bugs fixed during development (see
the ``why this test exists'' notes throughout). A runnable test suite backing
this plan lives in \code{tests/}: \textbf{57 automated tests, all passing as of
this writing} --- 55 unit/integration/CSV-system tests plus 2 that bring up the
real Docker/Postgres pipeline end to end (see the Test Execution Plan for how to
run them).

%==============================================================================
\section{Test Coverage}
%==============================================================================

Four testing types are used, matched to where risk actually concentrates in this
project, as summarized in Table~\ref{tab:coverage}.

\begin{longtable}{@{}>{\RaggedRight\arraybackslash}p{0.145\textwidth}
                     >{\RaggedRight\arraybackslash}p{0.40\textwidth}
                     >{\RaggedRight\arraybackslash}p{0.385\textwidth}@{}}
\caption{Testing types, matched to where risk concentrates in the pipeline.}
\label{tab:coverage}\\
\toprule
\hd{Type} & \hd{What it covers} & \hd{Why it matters here} \\
\midrule
\endfirsthead
\multicolumn{3}{@{}l}{\small\itshape Table \ref{tab:coverage} continued}\\[2pt]
\toprule
\hd{Type} & \hd{What it covers} & \hd{Why it matters here} \\
\midrule
\endhead
\midrule
\multicolumn{3}{r@{}}{\small\itshape continued on next page}\\
\endfoot
\bottomrule
\endlastfoot

\textbf{Unit} &
Individual functions in \code{etl/build_features.py},
\code{etl/combine_schools.py} (its two pure-Python helpers only --- see gap
below), \code{etl/build_rigor_classification.py},
\code{etl/build_modeling_dataset.py} &
This is where the real bugs have lived: sector misclassification, IB match-tier
gating, LEAID derivation, the CEEB fan-out --- all were logic errors inside
single functions, catchable in isolation. \\
\addlinespace

\textbf{Integration} &
How the CSV-driven pipeline stages compose
(\code{build_features.build()} \ra \code{build_modeling_dataset}'s freeze steps) &
Catches column-contract breaks between stages that unit tests on either side,
run alone, cannot see. \\
\addlinespace

\textbf{System} &
The five-script CSV pipeline run end to end against the real exported data
(\code{csv_exports/}), \textbf{and} the full Docker/Postgres pipeline
(\code{docker compose up}, \code{etl/run_all.py}) against a freshly started
database &
Confirms the whole thing actually runs, not just that each piece is individually
correct; this is also where the missing \code{scikit-learn} dependency (see
\S\ref{sec:env}) was caught. \\
\addlinespace

\textbf{User Acceptance (UAT)} &
Planned reviews of the two data dictionaries, \code{modeling_dataset.csv}, and
the rigor/clustering/benchmarking outputs, by the project client (Bob Henkins,
NU Undergraduate Admissions, and Adam) and, per the assignment's ``we may engage
the other team,'' a second student team &
The platform's entire purpose (per the report's Section 1) is serving admissions
officers standardized, auditable context --- correctness by our own tests is
necessary but not sufficient; the client has to find the output usable. \\
\end{longtable}

The Docker/Postgres pipeline was verified end to end on 2026-07-17 while this plan
was being finalized: \code{docker compose up -d db} followed by
\code{docker compose run --rm etl} completed successfully in \textbf{$\sim$106
seconds} against the raw data under \code{data/updated-sheng/}, and
\code{schools_org_all}'s CEEB match count of \textbf{16,508} matched the CSV-path
result exactly, cross-validating that both paths compute the same result from the
same join logic. That run is now covered by an automated test,
\code{tests/test_docker_pipeline.py}, which \emph{skips} rather than fails when
Docker isn't reachable or the raw data isn't present --- the same pattern
\code{test_system_pipeline.py} already used for the CSV-only path.

One gap remains, no longer undocumented but still not automatic. The
reproducibility gap for a genuinely fresh clone is real, and running the pipeline
once locally does not remove it. The directory \code{data/updated-sheng/} --- which
holds Sheng's combined schools export, the client's org export, and the raw
CRDC/EDFacts data --- is gitignored, and correctly so: it totals roughly
\textbf{2.6~GB} (a single CRDC folder alone is 794~MB, and two EDFacts files are
938~MB and 875~MB, each individually over GitHub's 100~MB per-file limit), far past
what is reasonable to version in git. As of 2026-07-17 this is a documented manual
step rather than an unwritten gap: the main \code{README.md} now links the team's
Google Drive folder holding the data and instructs a contributor to download it
into \code{data/updated-sheng/} before running the Docker path. A new team member
or a CI runner cloning the repo fresh must still perform that download manually ---
\code{test_docker_pipeline.py} will correctly skip until they do --- but the
hand-off itself is written down. Section~\ref{sec:env} gives the full setup
instructions, along with the remaining nice-to-have of a small sanitized fixture
that would let this path run in CI without the full 2.6~GB.

Relatedly, \code{etl/combine_schools.py}'s SQL-backed join functions
(\code{build_schools_org_enriched}, \code{build_schools_org_all},
\code{build_public_schools_enriched}, \code{build_private_schools_enriched}, and
\code{build_cps_nces_crosswalk}) are now exercised end to end by
\code{test_docker_pipeline.py} and its row-count cross-check, a real improvement
over having had zero coverage. This remains system-level assurance --- that the
whole run succeeded with the right row count --- rather than unit-level coverage of
each function's specific join semantics, so dedicated unit tests for these
functions, isolated from a live database, are still a gap.

Two areas are explicitly \textbf{not} covered by this plan, documented rather than
silently skipped. Load and performance testing is out of scope, because this is a
batch ETL pipeline re-run at most annually (per the report's reproducibility goal),
not a service under concurrent load. Security testing is likewise deferred: the
platform has no external-facing API or authentication surface yet, and this should
be revisited if one is added.

%==============================================================================
\section{Test Scenarios}
%==============================================================================

\subsection*{Unit test scenarios (by module)}

\noindent\textbf{\code{build_features.py}.} The unit tests cover bucket-midpoint
parsing across every observed format --- exact ranges (\code{"26\% - 50\%"}),
open-ended high (\code{"Over 90\%"}, \code{"greater than 20"}), open-ended low
(\code{"10\% or fewer"}), bare numbers (\code{"0"}), unparseable strings, and null
input. Winsorization must clip extreme outliers at the 1st/99th percentile without
crashing on all-null input. Sector classification is checked for a public HS, a
private HS via \code{nu_type}, a private HS via a school-side \code{pss_id} record
with \textbf{no} \code{nu_type} (the exact case a real bug missed), an org-only row
with no school-side match, and the mutual exclusivity of the public and private
flags. IB flag gating verifies that a \code{review}-tier match counts as a
candidate while a \code{reject}-tier match does not (a real bug counted both) and
that no-match rows are \code{0} rather than null. Finally, LEAID derivation must
take \code{nces_id_12[:7]} and return null when \code{nces_id_12} is itself null,
never falling back to the broken 5-character \code{leaid} column.

\noindent\textbf{\code{combine_schools.py}.} Name normalization maps
\code{"Saint"} \ra \code{"ST"} and \code{"HS"} \ra \code{"HIGH SCHOOL"}, strips
punctuation, and is null-safe. CEEB tie resolution keeps the CEEB on the best-tier
match while nulling and flagging the loser; ties at equal tier prefer an exact name
match; a frame with no duplicates is left untouched; and a unique CEEB is never
modified.

\noindent\textbf{\code{build_rigor_classification.py}.} Z-scoring must produce
mean-zero/std-one output on normal input, all-NaN output on zero-variance input
(never dividing by zero), and NaN passthrough. The weighted composite must match a
plain weighted average under full coverage, \textbf{reallocate weight
proportionally} when a component is missing (rather than zero-imputing it or
dropping the row), yield NaN (not 0) when zero components are available, and ignore
a zero-weight component --- IB in the default scheme --- even when it is present.
Tier assignment must produce five roughly-equal quintile buckets in the correct
ordinal order (lowest score \ra ``Below Average''), leaving NaN scores untiered
rather than defaulting them into a tier.

\noindent\textbf{\code{build_modeling_dataset.py}.} The minimum-size freeze drops
schools below the enrollment floor but \textbf{keeps} schools with unknown (null)
enrollment rather than treating null as ``too small,'' and the boundary value is
inclusive. Universe restriction leaves only public and private HS rows, dropping
\code{other/oos}. The sentinel scrub turns negative sentinel codes in rate- and
score-named columns into null while leaving a legitimate \code{0} untouched, and it
never scrubs a non-suspect column even when that column holds a value that would be
a sentinel elsewhere.

\subsection*{Integration test scenarios}

The synthetic fixture --- covering a public HS, a private HS via \code{pss_id}, a
private HS via \code{nu_type}, and an org-only row --- must survive
\code{build_features.build()} \ra universe restriction \ra the minimum-size freeze
with its sector labels intact and with no row collapsing into another. Chaining
these steps must also produce no duplicate school identities, a stand-in for the
class of bug the real CEEB fan-out was.

\subsection*{System test scenarios (automated, verified)}

Each of the five CSV pipeline scripts --- \code{build_features.py},
\code{build_modeling_dataset.py}, \code{build_rigor_classification.py},
\code{build_clustering.py}, and \code{build_benchmarking.py} --- must run to
completion against the real \code{csv_exports/} data and produce output of the
expected shape: a row-count floor met, the required columns present, and
categorical outputs constrained to their valid value sets. All five pass as of this
writing. Separately, \code{docker compose up} must bring up Postgres 16 and the ETL
container, \code{etl/run_all.py} must complete against it, and
\code{schools_org_all}'s row count must then match the CSV-path result; both pass,
completing in $\sim$106 seconds. That second test runs only when Docker is
reachable and the raw data is present locally, auto-skipping otherwise, since
neither condition holds on a fresh clone (see the gap below and
\S\ref{sec:env}).

\subsection*{Remaining gap --- not this pipeline's correctness, but who can run it}

The Docker/Postgres system test above is real and passing, but only on machines
that already have \code{data/updated-sheng/}'s $\sim$2.6~GB present locally
(correctly gitignored, not a bug --- see \S\ref{sec:env}). A fresh clone, or a CI
runner, still cannot run it without a manual download step first --- now documented
in \code{README.md} as a linked team Drive folder, not an unwritten gap, but still
a step a human must perform rather than something automated. A small sanitized
fixture of the data, scoped to just enough rows to exercise the joins, remains a
nice-to-have for eventually running this test in CI without the full 2.6~GB, and is
not built yet.

\subsection*{UAT scenarios (planned --- none of these have happened yet)}

Three acceptance checks are planned. First, Bob and Adam should be able to locate,
in \code{data_dictionary_schools_org_enriched.csv} and
\code{data_dictionary_modeling_dataset.csv}, the source, vintage, and description
for any variable they ask about --- literally what they requested in the
2026-07-14 meeting. Second, a reviewer unfamiliar with the pipeline should be able
to open \code{modeling_dataset_v1_2026-07-17.csv} and correctly interpret
\code{rigor_tier_label}, \code{sector}, and \code{funding_source} from the data
dictionary alone, without reading the ETL code. Third, the ``other team'' named in
the assignment's UAT engagement note should be able to reproduce one full pipeline
run from \code{csv_exports/} independently and obtain matching row counts, without
anyone on our team walking them through it live.

%==============================================================================
\section{Test Case Design}
%==============================================================================

Full test cases live as executable code in \code{tests/}, not merely described
here, so they cannot drift out of sync with the actual pipeline.
Table~\ref{tab:cases} lists representative examples in the format the rubric asks
for --- steps, data input, and expected result.

{\footnotesize
\begin{longtable}{@{}>{\RaggedRight\arraybackslash}p{0.055\textwidth}
                     >{\RaggedRight\arraybackslash}p{0.155\textwidth}
                     >{\RaggedRight\arraybackslash}p{0.25\textwidth}
                     >{\RaggedRight\arraybackslash}p{0.17\textwidth}
                     >{\RaggedRight\arraybackslash}p{0.25\textwidth}@{}}
\caption{Representative test cases (steps / data input / expected result).}
\label{tab:cases}\\
\toprule
\hd{ID} & \hd{Test} & \hd{Input} & \hd{Steps} & \hd{Expected result} \\
\midrule
\endfirsthead
\multicolumn{5}{@{}l}{\itshape Table \ref{tab:cases} continued}\\[2pt]
\toprule
\hd{ID} & \hd{Test} & \hd{Input} & \hd{Steps} & \hd{Expected result} \\
\midrule
\endhead
\midrule
\multicolumn{5}{r@{}}{\itshape continued on next page}\\
\endfoot
\bottomrule
\endlastfoot

UT-01 & \code{test_is_private_hs_via_pss_id_only} &
Fixture row: \code{school_id="S2"}, \code{pss_id="P1"}, \code{nu_type=NaN},
\code{school_level=NaN} & Run \code{build(df)} &
\code{is_private_hs == True}, \code{sector == "private"} --- regression test for a
real bug where such rows fell into \code{"other/oos"}. \\
\addlinespace

UT-02 & \code{test_reject_tier_does_not_count} &
Fixture row with \code{ib_school_id} set, \code{ib_match_tier="reject"} &
Run \code{build(df)} &
\code{ib_flag_candidate == 0} --- regression test for a real bug where reject-tier
matches were counted as confirmed. \\
\addlinespace

UT-03 & \code{test_missing_component_reallocates_weight_proportionally} &
\code{comp = \{"a":[1, NaN], "b":[3,4], "c":[5,6]\}},
\code{weights=\{"a":.5,"b":.25,"c":.25\}} &
Call \code{weighted_composite(comp, weights)} &
Row 1 score $=$ \code{4*.5 + 6*.5} (weight renormalized over available
components). \\
\addlinespace

UT-04 & \code{test_keeps_the_best_tier_match} &
Two schools sharing CEEB \code{050222}: one \code{auto_accept} exact-name match,
one \code{review}-tier different school &
Call \code{resolve_ceeb_ties(df)} &
\code{auto_accept} row keeps the CEEB; \code{review} row's CEEB is nulled and
flagged. \\
\addlinespace

IT-01 & \code{test_full_chain_on_fixture} &
4-row fixture (public HS, 2 private HS variants, 1 org-only) &
\code{build()} \ra \code{restrict_to_hs_universe()} \ra
\code{apply_min_size_freeze()} &
Only public/private rows remain; each is public XOR private; the
\code{pss_id}-only private school survived the chain. \\
\addlinespace

ST-01 & \code{test_build_rigor_classification} (system) &
Real \code{modeling_dataset_v1_*.csv} (34,392 rows) &
Run \code{build_rigor_classification.py} via subprocess &
Exit code 0; output has \code{rigor_tier_label}; every non-null value is one of
the five valid tier names. \\
\addlinespace

ST-02 & \code{test_full_pipeline_completes_successfully}
(\code{tests/test_docker_pipeline.py}) &
Raw data under \code{data/updated-sheng/} --- present on this development machine,
not on a fresh clone ($\sim$2.6~GB, gitignored) &
Bring up Postgres via \code{docker compose up -d db}, wait healthy,
\code{docker compose run --rm etl} &
Exit code 0; ``Pipeline completed successfully'' in output. Auto-skips if Docker
or the raw data isn't present, so it's automated where its precondition holds
rather than manual everywhere. \\
\addlinespace

ST-03 & \code{test_schools_org_all_row_count_matches_csv_path}
(\code{tests/test_docker_pipeline.py}) &
Same as ST-02 &
Query \code{schools_org_all} row count from the live database after ST-02 &
Matches the CSV-path row count (53,966) --- cross-validates both paths
independently. \\
\end{longtable}
}

Adding a test case follows one convention: any new bug found gets a regression test
named for the bug, not just for the function it lived in. This matches the pattern
already used throughout, described directly in \code{test_build_features.py}'s
docstring.

%==============================================================================
\section{Test Environment Setup}\label{sec:env}
%==============================================================================

Two environments are supported, matching how the pipeline is actually run day to
day. The first, \textbf{CSV-only (no database)}, is what most contributors use: it
requires Python 3.11+ (matching \code{etl/Dockerfile}'s \code{python:3.11-slim};
developed and tested here on 3.13.5) together with the packages pinned in
\code{etl/requirements.txt}, and it operates directly on \code{csv_exports/*.csv}
--- this is what the system tests of \S1 run against. The second, the \textbf{full
Docker stack}, is defined by \code{docker-compose.yml}, which brings up Postgres 16
(\code{capstone-db}, port 5433\ra5432) and an ETL container (\code{capstone-etl},
built from \code{etl/Dockerfile}) that runs \code{etl/run_all.py} against it; it is
needed for \code{combine_schools.py}'s SQL-backed builds and for anyone re-deriving
\code{csv_exports/} from raw source data rather than reading the already-exported
CSVs.

The Docker stack was verified working end to end on 2026-07-17:
\code{docker compose up -d db}, then \code{docker compose run --rm etl}, completed
successfully in $\sim$106 seconds and produced a \code{schools_org_all} row count
matching the CSV path exactly. That is covered by the automated test
\code{tests/test_docker_pipeline.py}.

As of 2026-07-17 the data hand-off is documented, not merely recommended. The
directory \code{data/updated-sheng/} --- Sheng's combined schools export, Bob's org
export, and the raw CRDC/EDFacts assessment data --- is gitignored, and correctly
so: it is roughly \textbf{2.6~GB} (CRDC data alone is 794~MB, and two EDFacts files
are 938~MB and 875~MB, each individually over GitHub's 100~MB per-file limit), far
past what is reasonable to version in git. That limit was checked directly against
GitHub's 100~MB hard per-file cap before attempting anything, rather than
discovered through a failed push. The data is instead distributed through a team
Google Drive folder linked from the main \code{README.md}'s Quickstart section
(request access from Max if it is not visible). A new contributor or CI runner still
needs to download it before the Docker/Postgres path or
\code{tests/test_docker_pipeline.py} will do anything beyond auto-skip, but that is
now a documented three-minute setup step rather than an undocumented gap. A small
sanitized fixture, scoped to just enough rows to exercise the joins, remains a
nice-to-have for eventually running this path in CI without the full 2.6~GB, but is
not built yet.

Test-specific dependencies are kept in \code{requirements-test.txt} (pytest 9.1.1),
separate from \code{etl/requirements.txt}, since pytest is a development- and
test-time dependency rather than a pipeline runtime one. Writing this plan already
caught one real environment gap: \code{etl/requirements.txt} was missing
\code{scikit-learn}, which \code{build_clustering.py} requires, meaning the
documented Docker environment could not actually run that script. It was fixed as
part of this work (see \code{etl/requirements.txt}) and confirmed by the successful
Docker build above, not merely asserted.

Configuration is read by \code{etl/config.py} from the environment variables
\code{DB_HOST}, \code{DB_PORT}, \code{DB_NAME}, \code{DB_USER}, and \code{DB_PASS}
(defaults \code{localhost:5432/capstone/postgres}), overridden by
\code{docker-compose.yml}'s \code{capstone}/\code{capstone}/\code{capstone} when run
in Docker, while \code{NU_MASTER_PATH} and \code{SCHOOLS_CEEB_PATH} point at the raw
data files under \code{data/}. The full suite is run as follows:

\begin{lstlisting}[style=shell]
pip install -r requirements-test.txt -r etl/requirements.txt
pytest tests/ -v
\end{lstlisting}

\noindent The system tests (\code{test_system_pipeline.py} and
\code{test_docker_pipeline.py}) auto-skip when their respective preconditions ---
exported CSVs, or Docker together with the raw data --- are not present, so the
suite still runs the unit and integration tests on a checkout without either.

%==============================================================================
\section{Test Execution Plan}
%==============================================================================

Execution is tied to the project's own weekly schedule rather than an invented
cadence, and Table~\ref{tab:exec} is honest about what is and is not formally
assigned. No RACI document exists in this repo: \code{docs/BOB_BRIEFING.md}
references ``per the RACI, master database \& data pulls is Max/Qifan's
workstream,'' which is the closest documented scope to this pipeline's code, but no
RACI covering \emph{testing} specifically has been written. Where a name appears
below it is inferred from that one existing reference, not from a formal
assignment, and is flagged as its own action item rather than silently assumed.

{\small
\begin{longtable}{@{}>{\RaggedRight\arraybackslash}p{0.175\textwidth}
                     >{\RaggedRight\arraybackslash}p{0.15\textwidth}
                     >{\RaggedRight\arraybackslash}p{0.315\textwidth}
                     >{\RaggedRight\arraybackslash}p{0.24\textwidth}@{}}
\caption{Test execution phases, timing, and (candidate) ownership.}
\label{tab:exec}\\
\toprule
\hd{Phase} & \hd{When} & \hd{What runs} & \hd{Who} \\
\midrule
\endfirsthead
\multicolumn{4}{@{}l}{\itshape Table \ref{tab:exec} continued}\\[2pt]
\toprule
\hd{Phase} & \hd{When} & \hd{What runs} & \hd{Who} \\
\midrule
\endhead
\midrule
\multicolumn{4}{r@{}}{\itshape continued on next page}\\
\endfoot
\bottomrule
\endlastfoot

Unit + integration, every commit &
Ongoing, starting now (Week 4) &
\code{pytest tests/test_build_features.py tests/test_combine_schools.py
tests/test_build_rigor_classification.py tests/test_build_modeling_dataset.py
tests/test_integration_pipeline.py} &
Whoever touches \code{etl/*.py} --- run locally before pushing; no CI is wired up
yet (see Defect Management), so this is currently a manual discipline, not an
enforced gate. \\
\addlinespace

System test, before every new versioned dataset &
Whenever \code{modeling_dataset_v<N>} is cut &
\code{pytest tests/test_system_pipeline.py} &
Whoever cuts the version. \\
\addlinespace

Full Docker system test &
Before any DB-schema-affecting change ships (new table, changed join); verified
passing on 2026-07-17 &
\code{pytest tests/test_docker_pipeline.py} (\code{docker compose up},
\code{etl/run_all.py}, then a row-count cross-check) --- auto-skips on machines
without Docker/the raw data (see \S\ref{sec:env}) &
Closest existing RACI scope: Max/Qifan (``master database \& data pulls'' per
\code{BOB_BRIEFING.md}) --- needs explicit confirmation, not assumed. \\
\addlinespace

Hyperparameter\slash{}validation-specific testing &
Week 7 (per the project's own schedule --- Section 4.6 of the report) &
Train/test/validation split checks, overfitting/underfitting checks --- not yet
built; this plan's system/integration tests are the pre-requisite groundwork for
that phase &
Not assigned --- no RACI line covers modeling/validation specifically; needs a
decision, not a default. \\
\addlinespace

UAT round 1 &
As soon as Bob/Adam have bandwidth (not gated on a specific week) &
Structured review of both data dictionaries $+$ \code{modeling_dataset.csv}, using
the UAT scenarios in \S2 &
Bob, Adam. \\
\addlinespace

UAT round 2 (peer team) &
Before Week 9 presentation prep &
The engaged ``other team'' independently reproduces one pipeline run and reviews
the rigor/clustering/benchmarking docs for clarity &
Second student team $+$ our team. \\
\end{longtable}
}

Several dependencies constrain this schedule. The system tests depend on
\code{csv_exports/} being current, so it is regenerated via the chain in \S1 before
they run. UAT round 2 depends on round 1 feedback having been incorporated first,
so the peer team is not reviewing something already known to be wrong. The Docker
system test runs today on any machine with Docker and the raw data present; a fresh
clone or CI runner needs the manual download step documented in \code{README.md}
first (see \S\ref{sec:env}), and stays unavailable until a sanitized fixture
removes that manual step entirely. Writing this plan also surfaced a standing action
item: a real RACI covering test ownership specifically does not exist yet, and the
team should produce one --- even a short one --- rather than leaving this plan to
infer ownership from a single line in \code{BOB_BRIEFING.md}.

%==============================================================================
\section{Defect Management}
%==============================================================================

No formal issue tracker is wired up yet, so this section states the process going
forward rather than one already running.

Defects found through testing, whether automated or via manual and UAT review, are
filed as GitHub Issues on \code{mchalekson/498-capstone}, tagged by the pipeline
stage affected (\code{etl}, \code{rigor-classification}, \code{clustering},
\code{benchmarking}, or \code{data-quality}). Each issue states what broke, the
failing test if the defect is automated or the specific input that triggered it
otherwise, and the expected versus actual output --- the same information already
present in this session's fix commits, of which the CEEB fan-out fix is an example,
documenting the exact failing pattern with real school names and CEEB values.

Severity is assigned on three tiers. A \emph{blocking} defect breaks a pipeline
stage outright, as the missing \code{scikit-learn} dependency did, and is fixed
before the next system-test run. A \emph{data-quality} defect produces output that
is wrong --- the CEEB fan-out and the IB sector-classification bug are examples ---
and is fixed before the next versioned \code{modeling_dataset} cut, and documented
in the relevant \code{docs/*.md} memo regardless of when the fix lands. A
\emph{documentation} defect is a case where code and docs disagree but the code
itself is correct, such as the \code{parse_bucket_midpoint} docstring typo found
while writing this plan's tests; it is low priority and batched into the next
doc-touching commit.

Issue status --- open, in progress, or resolved --- lives on the GitHub Issue
itself, with no separate spreadsheet, to avoid the two-sources-of-truth problem
this project has already flagged for data-vintage tracking. Every defect fix ships
with a regression test named for the bug (per the convention in \S3), so a defect
is never merely ``resolved,'' it is resolved-and-covered.

%==============================================================================
\section{Regression Testing}
%==============================================================================

The full unit-and-integration suite (\code{pytest tests/ -k "not system"}) runs
before every commit that touches \code{etl/*.py}, rather than on a fixed calendar
cadence --- appropriate for a small team iterating quickly rather than a large one
needing a scheduled batch window. Every previously found bug in this pipeline has a
corresponding test --- the sector-classification bug, the IB gating bug, the LEAID
derivation bug, the CEEB fan-out, and the docstring/behavior mismatch in
\code{parse_bucket_midpoint} --- so regression here specifically means ``did a fixed
bug come back,'' not just ``did anything change.''

At the full-pipeline level, the system-test suite reruns the entire five-script CSV
chain against the current \code{csv_exports/} data whenever any script in that chain
changes, catching regressions that only manifest at real production scale, where
small-fixture unit tests can miss scale-dependent issues such as the quintile-tier
bucketing needing a large-enough population. On machines with Docker and the raw
data available, \code{test_docker_pipeline.py} extends this to the database-backed
join functions as well, cross-checking the DB path's row counts against the CSV
path's --- something not yet possible in CI (see \S\ref{sec:env}). That absence of
CI is a real gap rather than one glossed over: the recommendation is to wire up
GitHub Actions to run the non-system tests on every push once the team has
bandwidth, so regression testing stops depending on individual discipline.

%==============================================================================
\section{Documentation and Reporting}
%==============================================================================

Test results are reported the same way every other finding in this project has
been --- as dated updates in the relevant \code{docs/*.md} file
(\code{EDA_features_joined.md}, \code{RIGOR_CLASSIFICATION.md},
\code{CLUSTERING.md}, \code{BENCHMARKING.md}), not as a separate test-report
artifact disconnected from the data findings it relates to. This test plan itself
(\code{docs/TEST_PLAN.md}) is the canonical reference for test \emph{process}, while
individual \emph{results} stay attached to the deliverable they validate.

Reporting frequency follows the work: results are documented after every
system-test run that precedes a new versioned dataset, and immediately whenever a
defect is found rather than in batches, so that (per the severity tiers above)
blocking and data-quality defects are written into the relevant memo right away.

The client-facing summary channel is \code{docs/BOB_BRIEFING.md}, already the
established pattern for this --- it has flagged Bob-relevant items such as the CEEB
corruption finding throughout the project --- so test and defect findings that
affect data Bob or Adam would review are added there specifically, not only to the
internal EDA memos. UAT feedback from Bob, Adam, and the peer team is logged in the
same \code{BOB_BRIEFING.md} update sections, matching its existing ``what we
resolved ourselves'' and ``what we still need from you'' structure, so that
acceptance feedback and technical defect tracking do not fragment into separate,
hard-to-reconcile systems.

%==============================================================================
\section{Sample Test Report}
%==============================================================================

Per the assignment instructions (``test reports\ldots with sample/test data,
inputs, and expected outputs''), the run below is real rather than illustrative ---
captured 2026-07-17 against this project's actual code and data, not a mocked-up
example.

\subsection*{Automated suite --- full run}

\begin{lstlisting}[style=shell]
$ pytest tests/ -v
...
57 passed in 115.31s (0:01:55)
\end{lstlisting}

\noindent All 57 tests across all six test files passed:
\code{test_build_features.py} (21), \code{test_build_modeling_dataset.py} (7),
\code{test_build_rigor_classification.py} (11), \code{test_combine_schools.py}
(9), \code{test_docker_pipeline.py} (2), \code{test_integration_pipeline.py} (2),
and \code{test_system_pipeline.py} (5). This run included the Docker-gated tests,
meaning the raw source data was present and the full Postgres pipeline was
exercised rather than skipped.

\subsection*{Sample test case, executed --- UT-01 (private-school sector classification)}

\begin{table}[h]
\centering
\begin{tabularx}{\textwidth}{@{}>{\bfseries\RaggedRight\arraybackslash}p{0.22\textwidth}
                                >{\RaggedRight\arraybackslash}X@{}}
\toprule
\normalfont\textbf{Field} & \textbf{Value} \\
\midrule
Test data / input & A fixture row: \code{school_id="S2"}, \code{pss_id="P1"},
\code{nu_type=NaN}, \code{school_level=NaN} (see \code{tests/conftest.py}). \\
\addlinespace
Action & \code{build(tiny_schools_org_all)}. \\
\addlinespace
Expected output & \code{is_private_hs == True}, \code{sector == "private"}. \\
\addlinespace
Actual output & \code{is_private_hs == True}, \code{sector == "private"}. \\
\addlinespace
Result & \textbf{PASS}. \\
\bottomrule
\end{tabularx}
\end{table}

\subsection*{Sample test case, executed --- ST-02/ST-03 (Docker/Postgres system test)}

\begin{table}[h]
\centering
\begin{tabularx}{\textwidth}{@{}>{\bfseries\RaggedRight\arraybackslash}p{0.22\textwidth}
                                >{\RaggedRight\arraybackslash}X@{}}
\toprule
\normalfont\textbf{Field} & \textbf{Value} \\
\midrule
Test data / input & Raw source data under \code{data/updated-sheng/}
($\sim$2.6~GB --- see \S\ref{sec:env}); a freshly created Postgres 16 container. \\
\addlinespace
Action & \code{docker compose up -d db}, wait healthy,
\code{docker compose build etl}, \code{docker compose run --rm etl}, then query
\code{schools_org_all} row count. \\
\addlinespace
Expected output & Exit code 0; ``Pipeline completed successfully'' in output;
\code{schools_org_all} row count matches the CSV-path result (53,966). \\
\addlinespace
Actual output & Exit code 0; pipeline completed in $\sim$106 seconds;
\code{schools_org_all} row count $=$ 53,966, exactly matching the CSV path. \\
\addlinespace
Result & \textbf{PASS}. \\
\bottomrule
\end{tabularx}
\end{table}

\subsection*{Sample query against the live database --- real data, not a mock}

To demonstrate that the database is genuinely queryable end to end, two arbitrary
example queries were run against \code{schools_org_all} (53,966 rows) in the running
Postgres container. The first returns the schools with the highest AP-test-taking
intensity:

\begin{lstlisting}[style=sql]
SELECT school_name, state, nu_avg_num_ap_tests_taken, nu_avg_freshman_sat
FROM schools_org_all
WHERE nu_avg_num_ap_tests_taken IS NOT NULL
ORDER BY nu_avg_num_ap_tests_taken DESC
LIMIT 5;
\end{lstlisting}

\begin{table}[h]
\centering
{\small\setlength{\tabcolsep}{5pt}
\begin{tabular}{@{}llrr@{}}
\toprule
\hd{School} & \hd{State} & \hd{Avg AP tests\slash{}student} & \hd{Avg freshman SAT} \\
\midrule
School of Science and Engineering   & TX & 15.49 & 1290 \\
School for the Talented and Gifted  & TX & 14.16 & 1250 \\
BASIS Chandler                      & AZ & 12.33 & 1260 \\
BASIS Peoria                        & AZ & 11.76 & 1200 \\
IDEA College Preparatory Pharr      & TX & 11.58 & 1080 \\
\bottomrule
\end{tabular}}
\end{table}

The second aggregates AP course counts and graduation rates by state, restricted to
public high schools in states with 500 or more schools:

\begin{lstlisting}[style=sql]
SELECT state, COUNT(*) AS n_schools,
       ROUND(AVG(crdc_ap_courses)::numeric, 1) AS avg_ap_courses_offered,
       ROUND(AVG(grad_rate_2021)::numeric, 1) AS avg_grad_rate
FROM schools_org_all
WHERE school_level = 'High'
GROUP BY state HAVING COUNT(*) > 500
ORDER BY avg_grad_rate DESC LIMIT 5;
\end{lstlisting}

\begin{table}[h]
\centering
\begin{tabular}{@{}lrrr@{}}
\toprule
\hd{State} & \hd{\# Schools} & \hd{Avg AP courses offered} & \hd{Avg grad rate} \\
\midrule
TX & 1,959 & 12.5 & 88.8\% \\
PA & 716   & 11.0 & 88.5\% \\
WI & 560   & 10.0 & 87.9\% \\
MO & 658   & 9.8  & 87.6\% \\
NY & 1,182 & 10.9 & 86.5\% \\
\bottomrule
\end{tabular}
\end{table}

Worth noting honestly rather than cropping out of the sample: this same query
returns a null average grad rate for Washington state (657 schools). That is not a
bug but a real coverage gap already documented in the data dictionary, since EDFacts
grad-rate data is not complete for every state. Including it here is deliberate,
because a sample test report that shows only clean rows is not an honest sample.

\end{document}
